% Volume I: The Epistemology of Suspicion
% Part II. Events, Records, and Reconstructions
% Chapter 10: Reconstruction as an Inverse Problem
% Spine: proof-spines/ch010-spine.md

\chapter{Reconstruction as an Inverse Problem}
\label{ch:ch010}


Historical and legal reconstruction are inverse problems under partial observation. Given evidence \(E\), one seeks
\[
\widehat H
=
\arg\max_{H\in\mathcal H}
P(H\mid E).
\]
But that posterior never arises from evidence alone. It depends on the hypothesis space \(\mathcal H\), the prior distribution over candidate histories, the likelihood model connecting traces to causes, and the exclusion or omission of unknown alternatives.

Bayesian language therefore does not eliminate judgment; it relocates judgment into model construction. Two investigators can agree on the surviving record and still disagree about priors, admissible hypotheses, or which absences count as informative. The chapter closes Part II by treating reconstruction as disciplined inference under irreversible loss: rigor comes not from pretending the inverse problem disappears, but from making the structure of the inference explicit.
